\label{sec:conclusion}

The relationship between exchange rate volatility and economic growth remains unsettled and also the question of how volatility should be measured to find its relationship with growth.
This study add a case instead of an answer.
Exchange rate volatility against the euro shows no independent relationship with growth in the six European Union members that kept their own currencies.
Where a relationship appears it runs trade exposure to the euro area and only in the years of financial crisis.

\vspace{\baselineskip}

Exchange rate volatility does not enter the model independently.
In the country and years fixed effects specification the volatility and its interaction with trade intensity are jointly significant.
There is no evidence that volatility affects growth at an average level of euro area trade exposure.
So the trade exposure is important. 
The joint test of trade intensity and the interaction rejects and the interaction coefficient is negative.
The marginal effect of volatility decreases steadily across the range of trade exposure from positive where trade with the euro are is small share of GDP to negative where it approaches the whole part.
This is the direction of the theory predicts where a country trade little with the euro area a moving currency absorbs shocks.
When it trades increase the cost that volatility imposes on trade cancels that benefit out.
None of these estimates is statistically significant and the confidence interval covers zero at every point.
Our analysis result holds only during financial crisis.
Excluding the crisis years removes the significance.
Excluding the COVID years leaves the estimate unchanged.
The analysis shows it is specific to 2008 to 2009 instead of periods of financial stress in general.

\vspace{\baselineskip}

These findings is different from Morina et al. (\citeyear{morina_effect_2020}) who present a significant negative effect of volatility on growth in a panel of fourteen Central and Eastern European countries.
There are three differences in design can be the reason.
The composition of the sample might be one of the reasons.
Their period covers years in which several members adopted the euro and they present that adoption was itself negative for growth.
Restricting the sample to countries that retained their currencies remove the source of variation.

\vspace{\baselineskip}

Also the measurement of volatility is different.
A real effective exchange rate shows movement against a country's trading partners.
And a euro exchange rate shows movement against the euro and nothing else.
These two do not behave the same and particularly for economics whose trade is diversified beyond the euro area.
The third difference is the model and it is the most important one.
Estimated only with country fixed effect finds a negative and significant coefficient on volatility.
This result changed after adding year affects which absorb what happened to all six countries in 2009 and would read as an effect of volatility.
It is not possible to say here which of the three is important.
The estimate is sensitive to the treatment of common shocks deserves attention.
The trade mechanism was also the least conclusive of the three from Morina et al. (\citeyear{morina_effect_2020}) study.
In both studies the interaction between volatility and trade exposure is harder to identify the direct effects.

\vspace{\baselineskip}

In general there are three directions can be address for further research.
The instrument for exchange rate volatility and following Janus and Riera-Crichton (\citeyear{thorsten_pdf_2015}) who use volatility of commodity terms of trade.
And a larger panel that extending to European countries outside the Union would fill the gap in the trade intensity range and reduce the sensitivity of the estimate to individual countries.
And estimating the volatility model and the panel regression together by bootstrap would produce standard errors that reflect the uncertainty in the first step which this analysis does not account for.

